\subsection{Related Work}

A limited number of previous works have learned dictionaries of sparsely-activating features in pre-trained models, including \citet{yun2021transformer} and \citet{sharkey2023technical}, the latter of which motivated this work. However, similar methods have been applied in other domains, in particular in understanding neurons in the visual cortex \citep{olshausen2004sparse, wright2022high}. 

In contrast to our approach, where we try to impose sparsity after training, many previous works have encouraged sparsity in neural networks via changes to the architecture or training process. These approaches include altering the attention mechanism \citep{correia2019adaptively}, adding $\ell^1$ penalties to neuron activations \citep{kasioumis2021elite, georgiadis2019accelerating}, pruning neurons \citep{frankle2018lottery}, and using the softmax function as the non-linearity in the MLP layers \citep{elhage2022solu}. However, training a state-of-the-art foundation model with these additional constraints is difficult \citep{elhage2022solu}, and improvements to interpretability are not always realized \citep{Meister2021IsSA}.

\subsection{Limitations and Future Work}

While we have presented evidence that our dictionary features are interpretable and causally important, we do not achieve 0 reconstruction loss (Equation \ref{eqn:loss_function}), indicating that our dictionaries fail to capture all the information in a layer's activations. We have also confirmed this by measuring the perplexity of the model's predictions when a layer is substituted with its reconstruction. For instance, replacing the residual stream activations in layer 2 of Pythia-70M with our reconstruction of those activations increases the perplexity on the Pile \citep{gao2020pile} from 25 to 40. To reduce this loss of information, we would like to explore other sparse autoencoder architectures and to try minimizing the change in model outputs when replacing the activations with our reconstructed vectors, rather than the reconstruction loss. Future efforts could also try to improve feature dictionary discovery by incorporating information about the weights of the model or dictionary features found in adjacent layers into the training process.

Our current methods for training sparse autoencoders are best suited to the residual stream. There is evidence that they may be applicable to the MLPs (see Appendix \ref{appendix:full_autointerp}), but the training pipeline used to train the dictionaries in this paper is not able to robustly learn overcomplete bases in the intermediate layers of the MLP. We're excited by future work investigating what changes can be made to better understand the computations performed by the attention heads and MLP layers, each of which poses different challenges.

In Section \ref{section:identifying_causal_features}, we show that for the IOI task, behaviour is dependent on a relatively small number of features. Because our dictionary is trained in a task-agnostic way, we expect this result to generalize to similar tasks and behaviours, but more work is needed to confirm this suspicion. If this property generalizes, we would have a set of features which allow for understanding many model behaviours using just a few features per behaviour. We would also like to trace the causal dependencies between features in different layers, with the overarching goal of providing a lens for viewing language models under which causal dependencies are sparse. This would hopefully be a step towards the eventual goal of building an end-to-end understanding of how a model computes its outputs.


\subsection{Conclusion}

Sparse autoencoders are a scalable, unsupervised approach to disentangling language model network features from superposition. Our approach requires only unlabelled model activations and uses orders of magnitude less compute than the training of the original models. We have demonstrated that the dictionary features we learn are more interpretable by autointerpretation, letting us pinpoint the features responsible for a given behaviour more finely, and are more monosemantic than comparable methods. This approach could facilitate the mapping of model circuits, targeted model editing, and a better understanding of model representations.

An ambitious dream in the field of interpretability is enumerative safety \citep{elhage2022toy}: producing a human-understandable explanation of a model's computations in terms of a complete list of the model's features and thereby providing a guarantee that the model will not perform dangerous behaviours such as deception. We hope that the techniques we presented in this paper also provide a step towards achieving this ambition.

